% Source: Wald GR, physical PDF p. 99 (printed p. 86).
% Coverage: Figure 4.3, asymptotically null hypersurfaces and radiated energy.
% Figures/tables/footnotes: Figure only; no tables or footnotes.
% Status: source-reviewed and compile-clean.
% Uncertainties: none.
\begingroup
\renewcommand{\thefigure}{4.3}
\begin{figure}[htbp]
  \centering
  \begin{tikzpicture}[scale=0.78,>=Stealth]
    % Timelike cylindrical boundary.
    \draw[thick] (-2.55,-1.50) arc[start angle=180,end angle=360,x radius=2.55,y radius=.62];
    \draw[thick,dashed] (2.55,-1.50) arc[start angle=0,end angle=180,x radius=2.55,y radius=.62];
    \draw[thick] (-2.55,2.35) arc[start angle=180,end angle=360,x radius=2.55,y radius=.62];
    \draw[thick] (2.55,2.35) arc[start angle=0,end angle=180,x radius=2.55,y radius=.62];
    \draw[thick] (-2.55,-1.50)--(-2.55,2.35);
    \draw[thick] (2.55,-1.50)--(2.55,2.35);
    \node[left] at (-2.58,.40) {\(S\)};

    % Asymptotically null hypersurfaces.
    \draw[thick] (-2.55,-1.50)
      .. controls (-1.55,-2.05) and (-.80,-2.32) .. (0,-2.33)
      .. controls (.80,-2.32) and (1.55,-2.05) .. (2.55,-1.50);
    \draw[thick] (-2.55,-1.50)
      .. controls (-1.55,-.92) and (-.80,-.72) .. (0,-.68)
      .. controls (.80,-.72) and (1.55,-.92) .. (2.55,-1.50);
    \draw[thick] (-2.55,2.35)
      .. controls (-1.55,1.77) and (-.80,1.56) .. (0,1.53)
      .. controls (.80,1.56) and (1.55,1.77) .. (2.55,2.35);
    \node[right] at (2.58,-1.40) {\(\Sigma_1\)};
    \node[right] at (2.58,2.28) {\(\Sigma_2\)};

    % Radiation reaches the lower null surface and timelike boundary.
    \foreach \x/\y/\dx/\dy in {
      -1.35/-3.05/.70/1.55,
      1.25/-2.95/-.60/1.48,
      -2.20/-.15/.55/-1.00,
      -.90/.25/.65/-1.15,
      .55/.20/-.25/-1.05,
      2.10/-.05/-.55/-1.00} {
      \draw[thick,decorate,decoration={snake,amplitude=2.6pt,segment length=9pt}]
        (\x,\y) -- ++(\dx,\dy);
    }
  \end{tikzpicture}
  \caption{Asymptotically null hypersurfaces \(\Sigma_1\) and \(\Sigma_2\).
  The gravitational radiation ``registers'' on \(\Sigma_1\) but does not
  ``register'' on \(\Sigma_2\). The energy, \(\Delta E\), carried off to infinity
  by gravitational radiation is given by an integral over the timelike hypersurface
  \(S\).}
  \label{fig:4.3}
\end{figure}
\endgroup
